\subsection{李代数的正则元}

{[}回忆一下，从现在起假定 \(k\) 是无限的。{]}

设 g 是一个维数为 n 的李代数。~若 \(x \in g\) ，把 ad x 的特征多项式写成如下形式

\[
\det (\mathrm{T} - \operatorname{ad} x) = \sum_ {i = 0} ^ {n} a _ {i} (x) \mathrm{T} ^ {i}, \quad \text { with } a _ {i} (x) \in k.
\]

我们有 \(a_i(x) = (-1)^{n - i}\mathrm{Tr}\left(\bigwedge^{n - i}\operatorname {ad}x\right)\) ，参看~《代数》，第三章，§8，第 11 节。这表明 \(x\mapsto a_i(x)\) 是一个 \(n - i\) 次齐次多项式映射，从 \(\mathfrak{g}\) 到 \(k\) （《代数》，第四章，§5，第 9 节）。

注. 1) 若 \(\mathfrak{g} \neq \{0\}\) ，则 \(a_0 = 0\) ，由于 (ad \(x(x) = 0\) 对所有 \(x \in \mathfrak{g}\) 成立。

\begin{enumerate}
\def\labelenumi{\arabic{enumi})}
\setcounter{enumi}{1}
\tightlist
\item
  设 \(k'\) 是 \(k\) 的一个扩张。写 \(\det(\mathrm{T} - \operatorname{ad} x') = \sum_{i=0}^{n} a_i'(x') \mathrm{T}^i\) ，对 \(x' \in \mathfrak{g} \otimes_k k'\) 。那么 \(a_i'| \mathfrak{g} = a_i\) 对所有 \(i\) 成立。
\end{enumerate}

定义 2. g 的秩，记作 rk(g)，是使 \(a_{l} \neq 0\) 的最小整数 l。g 的元素 x 称为正则元，如果 \(a_{l}(x) \neq 0\) 。

对所有 \(x \in \mathfrak{g}\) 有 \(\operatorname{rk}(\mathfrak{g}) \leq \dim \mathfrak{g}^0(x)\) ，且等式成立当且仅当 \(x\) 正则。

正则元之集在 \(\mathfrak{g}\) 中按 Zariski 拓扑是稠密且开的（附录 I）。

例. 1) 若 g 幂零，则 rk(g) = dim g，且 g 的所有元素都正则。

\begin{enumerate}
\def\labelenumi{\arabic{enumi})}
\setcounter{enumi}{1}
\tightlist
\item
  设 \(\mathfrak{g} = \mathfrak{sl}(2,k)\) 。若 \(x = \begin{pmatrix} \gamma & \alpha \\ \beta & -\gamma \end{pmatrix} \in \mathfrak{g}\) ，简单计算给出
\end{enumerate}

\[
\det (\mathrm{T} - \operatorname{ad} x) = \mathrm{T} ^ {3} - 4 (\alpha \beta + \gamma^ {2}) \mathrm{T}.
\]

若 \(k\) 的特征数 \(\neq 2\) ，则 \(\mathrm{rk}(\mathfrak{g}) = 1\) ，且正则元是那些 \(x\) ，使 \(\alpha \beta + \gamma^2 \neq 0\) 。

\begin{enumerate}
\def\labelenumi{\arabic{enumi})}
\setcounter{enumi}{2}
\tightlist
\item
  设 V 是一个有限维为 n 的向量空间，\(\mathfrak{g} = \mathfrak{gl}(V)\) 。设 \(\overline{k}\) 是 k 的一个代数闭包。设 \(x \in g\) ，并设 \(\lambda_{1}, \ldots, \lambda_{n}\) 是 x 的特征多项式在 \(\overline{k}\) 中的根（每个根写出的次数等于其重数）。从 \(V^{*} \otimes V\) 到 g 的典则同构与这两个空间的 g-模结构相容，换言之，它把 \(1 \otimes x - {}^{t}x \otimes 1\) 映到 ad x（第一章，§3，第 3 节，命题 4）。注意到 §1 命题 4 (i)，由此推出 ad x 的特征多项式的根是 \(\lambda_{i} - \lambda_{j}\) ，其中 \(1 \leq i \leq n, 1 \leq j \leq n\) （每个根写出的次数等于其重数）。于是，g 的秩是 n，且 x 正则当且仅当每个 \(\lambda_{i}\) 是 x 的特征多项式的单根。
\end{enumerate}

命题 6. 设 \(\mathfrak{g}\) 是一个李代数，\(k'\) 是 \(k\) 的一个扩张，\(\mathfrak{g}' = \mathfrak{g} \otimes_k k'\) 。

\begin{enumerate}
\def\labelenumi{(\roman{enumi})}
\tightlist
\item
  一个元素 \(x\) （属于 \(\mathfrak{g}\) ）在 \(\mathfrak{g}\) 中正则当且仅当 \(x \otimes 1\) 在 \(\mathfrak{g}'\) 中正则。\\
\item
  \(\operatorname{rk}(\mathfrak{g}) = \operatorname{rk}(\mathfrak{g}')\) 。
\end{enumerate}

这由注 2 得出。

命题 7. 设 \((\mathfrak{g}_i)_{i\in \mathrm{I}}\) 是一个有限的李代数族，\(\mathfrak{g} = \prod_{i\in \mathrm{I}}\mathfrak{g}_i\) 。

\begin{enumerate}
\def\labelenumi{(\roman{enumi})}
\tightlist
\item
  一个元素 \((x_{i})_{i\in \mathrm{I}}\) （属于 \(\mathfrak{g}\) ）在 \(\mathfrak{g}\) 中正则当且仅当对所有 \(i\in \mathrm{I}\) ，\(x_{i}\) 在 \(\mathfrak{g}_i\) 中正则。
\end{enumerate}

\[
\text {(ii)} \operatorname{rk} (\mathfrak {g}) = \sum_ {i \in I} \operatorname{rk} (\mathfrak {g} _ {i}).
\]

事实上，对任意 \(x = (x_{i})_{i \in \mathrm{I}} \in \mathfrak{g}\) ，\(ad_{g}x\) 的特征多项式是诸 \(ad_{g_{i}}x_{i}\) 的特征多项式的乘积。

命题 8. 设 \(f: \mathfrak{g} \to \mathfrak{g}'\) 是李代数的一个满射同态。

\begin{enumerate}
\def\labelenumi{(\roman{enumi})}
\item
  若 \(x\) 是 \(\mathfrak{g}\) 的正则元，则 \(f(x)\) 在 \(\mathfrak{g}'\) 中正则。若 \(\operatorname{Ker} f\) 包含在 \(\mathfrak{g}\) 的中心中，则逆命题成立。
\item
  \(\operatorname{rk}(\mathfrak{g}) \geq \operatorname{rk}(\mathfrak{g}')\) 。
\end{enumerate}

设 \(\mathrm{rk}(\mathfrak{g}) = r, \mathrm{rk}(\mathfrak{g}') = r'\) 。设 \(x \in \mathfrak{g}\) 。\(\operatorname{ad} x\) 、\(\operatorname{ad} f(x)\) 和 \(\operatorname{ad} x|\operatorname{Ker} f\) 的特征多项式形如

\[
\mathrm{P} (\mathrm{T}) = \mathrm{T} ^ {n} + a _ {n - 1} (x) \mathrm{T} ^ {n - 1} + \dots + a _ {r} (x) \mathrm{T} ^ {r},
\]

\[
\mathrm{Q} (\mathrm{T}) = \mathrm{T} ^ {n ^ {\prime}} + b _ {n ^ {\prime} - 1} (x) \mathrm{T} ^ {n ^ {\prime} - 1} + \dots + b _ {r ^ {\prime}} (x) \mathrm{T} ^ {r ^ {\prime}},
\]

\[
\mathrm{R} (\mathrm{T}) = \mathrm{T} ^ {n ^ {\prime \prime}} + c _ {n ^ {\prime \prime} - 1} (x) \mathrm{T} ^ {n ^ {\prime \prime} - 1} + \dots + c _ {r ^ {\prime \prime}} (x) \mathrm{T} ^ {r ^ {\prime \prime}},
\]

其中 \(a_i, b_i, c_i\) 是 \(\mathfrak{g}\) 上的多项式函数，且 \(a_r \neq 0\) ，\(b_{r'} \neq 0\) ，\(c_{r''} \neq 0\) 。我们有 \(\mathrm{P} = \mathrm{QR}\) ，故 \(r = r' + r''\) 且 \(a_r(x) = b_{r'}(x)c_{r''}(x)\) ，这就证明了 (ii) 以及 (i) 的第一个断言。若 \(\operatorname{Ker} f\) 包含在 \(\mathfrak{g}\) 的中心中，则 \(\mathrm{R(T)} = \mathrm{T}^{n''}\) ，从而 \(a_r(x) = b_{r'}(x)\) ，于是得到 (i) 的第二个断言。

推论. 设 \(\mathcal{C}_n\mathfrak{g}\) ( \(n \geq 0\) ) 是 \(\mathfrak{g}\) 的上升中心列的一项（第一章，§1，第 6 节）。\(\mathfrak{g}\) 的正则元是那些在 \(\mathfrak{g}/\mathcal{C}_n\mathfrak{g}\) 中的像正则的元素。

命题 9. 设 \(\mathfrak{g}\) 是一个李代数，\(\mathfrak{g}'\) 是 \(\mathfrak{g}\) 的一个子代数。\(\mathfrak{g}'\) 的每个在 \(\mathfrak{g}\) 中正则的元素在 \(\mathfrak{g}'\) 中正则。

对 \(x \in g'\) ，\(ad_{g}x\) 在 \(g'\) 上的限制是 \(ad_{g'}x\) ，从而通过转到商定义自同态 \(u(x)\) 于向量空间 \(g/g'\) 上。设 \(d_{0}(x)\) （相应地 \(d_{1}(x)\) ）是 \(ad_{g'}(x)\) （相应地 \(u(x)\) ）的幂零空间的维数，并设 \(c_{0}\) （相应地 \(c_{1}\) ）是 \(d_{0}(x)\) （相应地 \(d_{1}(x)\) ）当 x 属于 \(g'\) 时的最小值。存在非零多项式映射 \(p_{0}, p_{1}\) ，从 \(g'\) 到 k，使得

\[
d _ {0} (x) = c _ {0} \Longleftrightarrow p _ {0} (x) \neq 0, \quad d _ {1} (x) = c _ {1} \Longleftrightarrow p _ {1} (x) \neq 0.
\]

由于 k 无限，\(x \in g'\) 中使 \(d_{0}(x) = c_{0}\) 且 \(d_{1}(x) = c_{1}\) 的元素所成的集 S 非空。S 的每个元素在 \(g'\) 中正则。另一方面，S 是 \(g'\) 中使 \(ad_{g}x\) 的幂零空间有最小维数的元素之集，因而包含 \(g'\) 的每个在 g 中正则的元素。

注. 3) \(\mathfrak{g}'\) 的在 \(\mathfrak{g}\) 中正则的元素不一定存在。若至少存在一个，则这些元素之集正是上面证明中记作 S 的集合。

